% Part: proof-theory
% Chapter: normalization
% Section: permutations

\documentclass[../../../include/open-logic-section]{subfiles}

\begin{document}

\olfileid{pt}{nor}{per}

\olsection{د ځاے بدلولو بدلونونه}

په \Log{N1i} کښې پرېکونونه هغه برخې دي چې د
!!a{elimination} قاعدې په لويه مقدمه پاے ته رسېږي۔ د عادي کولو
په ثبوت کښې يو حالت دا دے چې د تر ټولو لوړې رتبې د پرېکونونو
او ورسره د !!{proof} د پرېکون اوږدوالے کم کړو۔ د $>1$ اوږدوالي
پرېکون په بېلو طريقو پاے ته رسېدے شي: دا په دې اړه لري چې په
پرېکون کښې د~$!A$ وروستے !!{formula} مورد د کومې قاعدې
(\Elim{\lor} يا \Elim{\lexists}) پايله ده، او د کوم
!!{elimination} استنتاج لويه مقدمه ده۔

مثلاً، که اړوند استنتاجونه \Elim{\lor} او \Elim{\land}
وي، $!A \ident !B \land !C$ وي او داسې فرعي
!!{proof}~$\delta_1$ وي چې په لاندې رسم پاے ته رسېږي:
\[
\AxiomC{}
\RightLabel{$\delta_2$}
\DeduceC{$!D \lor !E$}
\AxiomC{$\Discharge{!D}{x}$}
\RightLabel{$\delta_3$}
\DeduceC{$!B \land !C$}
\AxiomC{$\Discharge{!E}{y}$}
\RightLabel{$\delta_4$}
\DeduceC{$!B \land !C$}
\DischargeRule{\Elim{\lor}}{x\,y}
\TrinaryInfC{$!B \land !C$}
\RightLabel{\Elim{\land}[1]}
\UnaryInfC{$!B$}
\DisplayProof
\]
\textbf{د سرچينه‌يي پېژندنې سمون OLCMP-171:} دلته د پرېکون
فارمول د او مرکب دے، نۀ د يا مرکب؛ د لويې مقدمې پېژندنه د دواړو
اصلي رسمونو سره برابره شوه۔ خپله رسمونه عين ساتل شوي دي۔
په پرېکون کښې وروستے !!{formula} مورد~$!A_n$ د~\Elim{\lor}
پايله، يعنې لاندې $!B \land !C$ دے۔ له هغې يو مخکې
!!{formula} مورد~$!A_{n-1}$ د~\Elim{\lor} له کوچنيو مقدمو
څخه يوه ده۔

د دې پرېکون د اوږدوالي کمولو دپاره د \Elim{\lor} او
\Elim{\land} استنتاجونو ترتيب بدلولے، يعنې ``ځاے بدلولے'' شو۔
په~$\delta$ کښې فرعي ثبوت~$\delta_1$ په لاندې رسم بدلوو:
\[
\AxiomC{}
\RightLabel{$\delta_2$}
\DeduceC{$!D \lor !E$}
\AxiomC{$\Discharge{!D}{x}$}
\RightLabel{$\delta_3$}
\DeduceC{$!B \land !C$}
\RightLabel{\Elim{\land}[1]}
\UnaryInfC{$!B$}
\AxiomC{$\Discharge{!E}{y}$}
\RightLabel{$\delta_4$}
\DeduceC{$!B \land !C$}
\RightLabel{\Elim{\land}[1]}
\UnaryInfC{$!B$}
\DischargeRule{\Elim{\lor}}{x\,y}
\TrinaryInfC{$!B$}
\DisplayProof
\]
هغه پرېکونونه چې د \Elim{\lor} لاندې د \Elim{\land} په
مقدمه کښې ختمېدل، اوس د \Elim{\lor} د کوچنيو مقدمو د پاسه
د يوه \Elim{\land} استنتاج په مقدمه کښې ختمېږي۔ يعنې د هر
يو اوږدوالے په~$1$ کم شوے دے۔

سربېره پر دې، هر هغه پرېکون چې بشپړ په $\delta_2$،
$\delta_3$ يا $\delta_4$ کښې وي، يا په~$\delta$ کښې بشپړ له
$\delta_1$ څخه بهر وي، نۀ بدلېږي: هماغه !!{formula}s رانغاړي
او له هماغو استنتاجونو تېرېږي۔ نو په ځانګړي ډول د هغې رتبه
او اوږدوالے په~$\delta$ کښې د اړوند پرېکون سره عين دي۔ له دې
امله د نوي !!{proof} د پرېکون رتبه نۀ ده لوړه شوې او د نوي
!!{proof} د پرېکون اوږدوالے د زاړه !!{proof} څخه کم دے۔

\begin{prob}
د ځاے بدلولو له بدلون وروسته لږ تر لږه د يوه پرېکون اوږدوالے
په~$1$ کمېږي۔ که \Elim{\land} له يوه~\Elim{\lor} څخه تېر
کړو، په حقيقت کښې د دوو پرېکون برخو اوږدوالے کمېږي؛ نو په دې
حالت کښې د !!{proof} د پرېکون اوږدوالے لږ تر لږه په~$2$
کمېږي۔ آيا يو داسې بدلون د !!a{proof} د پرېکون اوږدوالے له~$2$
څخه زيات هم کمولے شي؟ آيا د پرېکون \emph{رتبه} هم ټيټېدې شي؟
\end{prob}

د \Elim{\lor} يا \Elim{\lexists} قاعدې د پاسه د
!!a{elimination} قاعدې د وړلو عمل د \emph{ځاے بدلولو بدلون}
بلل کېږي۔ د ځينو قاعده‌جوړو بڼې په
\olref{tab:permutation-conversions} کښې ورکړې دي۔ پاتې بڼې
د مشقونو په توګه پرېښودل شوې دي۔

% \begin{table}
% \begin{multline*}
% \shoveleft{\AxiomC{$!D \lor !E$}
% \AxiomC{$\Discharge{!D}{x}$}
% \shortDeduce
% \DeduceC{$!B \land !C$}
% \AxiomC{$\Discharge{!E}{y}$}
% \shortDeduce
% \DeduceC{$!B \land !C$}
% \DischargeRule{\Elim{\lor}}{x\,y}
% \TrinaryInfC{$!B \land !C$}
% \RightLabel{\Elim{\land}[1]}
% \UnaryInfC{$!B$}
% \DisplayProof}\\
% \shoveright{\longrightarrow\ 
% \AxiomC{$!D \lor !E$}
% \AxiomC{$\Discharge{!D}{x}$}
% \shortDeduce
% \DeduceC{$!B \land !C$}
% \RightLabel{\Elim{\land}[1]}
% \UnaryInfC{$!B$}
% \AxiomC{$\Discharge{!E}{y}$}
% \shortDeduce
% \DeduceC{$!B \land !C$}
% \RightLabel{\Elim{\land}[1]}
% \UnaryInfC{$!B$}
% \DischargeRule{\Elim{\lor}}{x\,y}
% \TrinaryInfC{$!B$}
% \DisplayProof}
% \\
% \shoveleft{\AxiomC{$!D \lor !E$}
% \AxiomC{$\Discharge{!D}{x}$}
% \shortDeduce
% \DeduceC{$!B \lif !C$}
% \AxiomC{$\Discharge{!E}{y}$}
% \shortDeduce
% \DeduceC{$!B \lif !C$}
% \DischargeRule{\Elim{\lor}}{x\,y}
% \TrinaryInfC{$!B \lif !C$}
% \AxiomC{$!B$}
% \RightLabel{\Elim{\lif}}
% \BinaryInfC{$!C$}
% \DisplayProof}\\
% \shoveright{\longrightarrow\ 
% \AxiomC{$!D \lor !E$}
% \AxiomC{$\Discharge{!D}{x}$}
% \shortDeduce
% \DeduceC{$!B \lif !C$}
% \AxiomC{$!B$}
% \RightLabel{\Elim{\lif}}
% \BinaryInfC{$!C$}
% \AxiomC{$\Discharge{!E}{y}$}
% \shortDeduce
% \DeduceC{$!B \lif !C$}
% \AxiomC{$!B$}
% \RightLabel{\Elim{\lif}}
% \BinaryInfC{$!C$}
% \DischargeRule{\Elim{\lor}}{x\,y}
% \TrinaryInfC{$!C$}
% \DisplayProof}
% \\
% \shoveleft{\AxiomC{$!D \lor !E$}
% \AxiomC{$\Discharge{!D}{x}$}
% \shortDeduce
% \DeduceC{$!B \lif !C$}
% \AxiomC{$\Discharge{!E}{y}$}
% \shortDeduce
% \DeduceC{$!B \lif !C$}
% \DischargeRule{\Elim{\lor}}{x\,y}
% \TrinaryInfC{$!B \lif !C$}
% \AxiomC{$!B$}
% \RightLabel{\Elim{\lif}}
% \BinaryInfC{$!C$}
% \DisplayProof}\\
% \shoveright{\longrightarrow\ 
% \AxiomC{$!D \lor !E$}
% \AxiomC{$\Discharge{!D}{x}$}
% \shortDeduce
% \DeduceC{$!B \lif !C$}
% \AxiomC{$!B$}
% \RightLabel{\Elim{\lif}}
% \BinaryInfC{$!C$}
% \AxiomC{$\Discharge{!E}{y}$}
% \shortDeduce
% \DeduceC{$!B \lif !C$}
% \AxiomC{$!B$}
% \RightLabel{\Elim{\lif}}
% \BinaryInfC{$!C$}
% \DischargeRule{\Elim{\lor}}{x\,y}
% \TrinaryInfC{$!C$}
% \DisplayProof}
% \\
% \shoveleft{\AxiomC{$!D \lor !E$}
% \AxiomC{$\Discharge{!D}{x}$}
% \shortDeduce
% \DeduceC{$\lforall[x][!B(x)]$}
% \AxiomC{$\Discharge{!E}{y}$}
% \shortDeduce
% \DeduceC{$\lforall[x][!B(x)]$}
% \DischargeRule{\Elim{\lor}}{x\,y}
% \TrinaryInfC{$\lforall[x][!B(x)]$}
% \RightLabel{\Elim{\lforall}}
% \UnaryInfC{$!B(t)$}
% \DisplayProof}\\
% \shoveright{\longrightarrow\ 
% \AxiomC{$!D \lor !E$}
% \AxiomC{$\Discharge{!D}{x}$}
% \shortDeduce
% \DeduceC{$\lforall[x][!B(x)]$}
% \RightLabel{\Elim{\lforall}}
% \UnaryInfC{$!B(t)$}
% \AxiomC{$\Discharge{!E}{y}$}
% \shortDeduce
% \DeduceC{$\lforall[x][!B(x)]$}
% \RightLabel{\Elim{\lforall}}
% \UnaryInfC{$!B(t)$}
% \DischargeRule{\Elim{\lor}}{x\,y}
% \TrinaryInfC{$!B(t)$}
% \DisplayProof}
% \end{multline*}
% \caption{Some permutation conversions for \Log{N1i}.}
% \ollabel{tab:permutation-conversions}
% \end{table}

{\small
\begin{longtable}{@{}l@{}}
\AxiomC{$!D \lor !E$}
\AxiomC{$\Discharge{!D}{x}$}
\shortDeduce
\DeduceC{$!B \land !C$}
\AxiomC{$\Discharge{!E}{y}$}
\shortDeduce
\DeduceC{$!B \land !C$}
\DischargeRule{\Elim{\lor}}{x\,y}
\TrinaryInfC{$!B \land !C$}
\RightLabel{\Elim{\land}[1]}
\UnaryInfC{$!B$}
\DisplayProof
$\longrightarrow$\\[6ex]
\quad\AxiomC{$!D \lor !E$}
\AxiomC{$\Discharge{!D}{x}$}
\shortDeduce
\DeduceC{$!B \land !C$}
\RightLabel{\Elim{\land}[1]}
\UnaryInfC{$!B$}
\AxiomC{$\Discharge{!E}{y}$}
\shortDeduce
\DeduceC{$!B \land !C$}
\RightLabel{\Elim{\land}[1]}
\UnaryInfC{$!B$}
\DischargeRule{\Elim{\lor}}{x\,y}
\TrinaryInfC{$!B$}
\DisplayProof
\\[10ex]
\AxiomC{$!D \lor !E$}
\AxiomC{$\Discharge{!D}{x}$}
\shortDeduce
\DeduceC{$!B \lif !C$}
\AxiomC{$\Discharge{!E}{y}$}
\shortDeduce
\DeduceC{$!B \lif !C$}
\DischargeRule{\Elim{\lor}}{x\,y}
\TrinaryInfC{$!B \lif !C$}
\AxiomC{$!B$}
\RightLabel{\Elim{\lif}}
\BinaryInfC{$!C$}
\DisplayProof
$\longrightarrow$\\[6ex]
\AxiomC{$!D \lor !E$}
\AxiomC{$\Discharge{!D}{x}$}
\shortDeduce
\DeduceC{$!B \lif !C$}
\AxiomC{$!B$}
\RightLabel{\Elim{\lif}}
\BinaryInfC{$!C$}
\AxiomC{$\Discharge{!E}{y}$}
\shortDeduce
\DeduceC{$!B \lif !C$}
\AxiomC{$!B$}
\RightLabel{\Elim{\lif}}
\BinaryInfC{$!C$}
\DischargeRule{\Elim{\lor}}{x\,y}
\TrinaryInfC{$!C$}
\DisplayProof
\\[10ex]
\AxiomC{$!D \lor !E$}
\AxiomC{$\Discharge{!D}{x}$}
\shortDeduce
\DeduceC{$!B \lif !C$}
\AxiomC{$\Discharge{!E}{y}$}
\shortDeduce
\DeduceC{$!B \lif !C$}
\DischargeRule{\Elim{\lor}}{x\,y}
\TrinaryInfC{$!B \lif !C$}
\AxiomC{$!B$}
\RightLabel{\Elim{\lif}}
\BinaryInfC{$!C$}
\DisplayProof
$\longrightarrow$\\[6ex]
\AxiomC{$!D \lor !E$}
\AxiomC{$\Discharge{!D}{x}$}
\shortDeduce
\DeduceC{$!B \lif !C$}
\AxiomC{$!B$}
\RightLabel{\Elim{\lif}}
\BinaryInfC{$!C$}
\AxiomC{$\Discharge{!E}{y}$}
\shortDeduce
\DeduceC{$!B \lif !C$}
\AxiomC{$!B$}
\RightLabel{\Elim{\lif}}
\BinaryInfC{$!C$}
\DischargeRule{\Elim{\lor}}{x\,y}
\TrinaryInfC{$!C$}
\DisplayProof
\\[10ex]
\AxiomC{$!D \lor !E$}
\AxiomC{$\Discharge{!D}{x}$}
\shortDeduce
\DeduceC{$\lforall[x][!B(x)]$}
\AxiomC{$\Discharge{!E}{y}$}
\shortDeduce
\DeduceC{$\lforall[x][!B(x)]$}
\DischargeRule{\Elim{\lor}}{x\,y}
\TrinaryInfC{$\lforall[x][!B(x)]$}
\RightLabel{\Elim{\lforall}}
\UnaryInfC{$!B(t)$}
\DisplayProof
\quad$\longrightarrow$\\[8ex]
\quad\AxiomC{$!D \lor !E$}
\AxiomC{$\Discharge{!D}{x}$}
\shortDeduce
\DeduceC{$\lforall[x][!B(x)]$}
\RightLabel{\Elim{\lforall}}
\UnaryInfC{$!B(t)$}
\AxiomC{$\Discharge{!E}{y}$}
\shortDeduce
\DeduceC{$\lforall[x][!B(x)]$}
\RightLabel{\Elim{\lforall}}
\UnaryInfC{$!B(t)$}
\DischargeRule{\Elim{\lor}}{x\,y}
\TrinaryInfC{$!B(t)$}
\DisplayProof
\\
\caption{د \Log{N1i} د ځاے بدلولو ځينې بدلونونه۔}
\ollabel{tab:permutation-conversions}
\end{longtable}
}

\begin{prob}
د \Elim{\lor} او ورپسې \Elim{\lor} يا \Elim{\lnot} د ځاے
بدلولو بدلونونه ورکړئ۔
\end{prob}

\begin{prob}
د \Elim{\lexists} او ورپسې \Elim{\exists} د ځاے بدلولو
بدلونونه ورکړئ۔ ووايئ چې په وروستي حالت کښې ولې د ځانګړي
متغير هېڅ شرط نۀ ماتېږي۔
\end{prob}

بايد ځان ډاډه کړو چې د ځاے بدلولو د بدلون په تطبيق سره
د ځانګړي متغير کوم شرط نۀ ماتوو۔ دا رسم وګورئ:
\[
\AxiomC{}
\RightLabel{$\delta_2$}
\DeduceC{$!D \lor !E$}
\AxiomC{$\Discharge{!D}{x}$}
\RightLabel{$\delta_3$}
\DeduceC{$\lexists[x][!B(x)]$}
\AxiomC{$\Discharge{!E}{y}$}
\RightLabel{$\delta_4$}
\DeduceC{$\lexists[x][!B(x)]$}
\DischargeRule{\Elim{\lor}}{x\,y}
\TrinaryInfC{$\lexists[x][!B(x)]$}
\AxiomC{$\Discharge{!B(c)}{z}$}
\RightLabel{$\delta_5$}
\DeduceC{$!C$}
\DischargeRule{\Elim{\lexists}}{z}
\BinaryInfC{$!C$}
\DisplayProof
\]
که دا په لاندې رسم بدل کړو:
\[
\AxiomC{}
\RightLabel{$\delta_2$}
\DeduceC{$!D \lor !E$}
\AxiomC{$\Discharge{!D}{x}$}
\RightLabel{$\delta_3$}
\DeduceC{$\lexists[x][!B(x)]$}
\AxiomC{$\Discharge{!B(c)}{z}$}
\RightLabel{$\delta_5$}
\DeduceC{$!C$}
\DischargeRule{\Elim{\lexists}}{z}
\BinaryInfC{$!C$}
\AxiomC{$\Discharge{!E}{y}$}
\RightLabel{$\delta_4$}
\DeduceC{$\lexists[x][!B(x)]$}
\AxiomC{$\Discharge{!B(c)}{z}$}
\RightLabel{$\delta_5$}
\DeduceC{$!C$}
\DischargeRule{\Elim{\lexists}}{z}
\BinaryInfC{$!C$}
\DischargeRule{\Elim{\lor}}{x\,y}
\TrinaryInfC{$!C$}
\DisplayProof
\]
اندېښنه کېدے شي چې ځانګړے متغير~$c$ مثلاً په~$!D$ کښې
راشي۔ ځکه فرض $!D$ د اصلي~$\Elim{\lexists}$ د پاسه ايستل
کېږي؛ نو په اصلي !!{proof} کښې دا د~\Elim{\exists} د لويې
مقدمې د پاسه په کوم پرانيستي فرض کښې نۀ راځي او د ځانګړي
متغير شرط هلته پوره دے۔ خو په نوي !!{proof} کښې $!D$~تر
هغې نۀ ايستل کېږي چې د~\Elim{\lexists} استنتاجونه تېر شي؛
نو د ځانګړي متغير شرط به مات شي۔ خو په ياد ولرئ چې زموږ
!!{proof}s منظم فرض شوي دي: هر ځانګړے متغير يوازې د يوه
استنتاج ځانګړے متغير دے او په !!{proof} کښې يوازې د
\Elim{\lexists} د کوچنۍ مقدمې د پاسه راځي۔ په ځانګړي ډول،
$c$~په $\delta_3$ يا~$\delta_4$ کښې نۀ شي راتلے۔

دا مثال بله خبره هم ښيي: په نوي !!{proof} کښې
د~$\delta_5$ دوه نقلونه شته۔ نو که $\delta_5$ د تر ټولو لوړې
رتبې پرېکون ولري، د پرېکون اوږدوالے مو \emph{نۀ} دے کم
کړے! ځکه بايد يوازې هغه وخت د ځاے بدلولو بدلون تطبيق کړو
چې په $\delta_5$ کښې د تر ټولو لوړې رتبې پرېکون نۀ وي۔ دا
د تل لپاره د تر ټولو \emph{پاسني} لوړې رتبې پرېکون په
غوره کولو يقيني کوو۔ يعنې داسې د تر ټولو لوړې رتبې پرېکون
غوره کوو چې فرعي !!{proof}، دلته~$\delta_1$، داسې بل د تر ټولو
لوړې رتبې پرېکون نۀ لري چې د همدغه فرعي !!{proof} د وروستي
استنتاج د پاسه ختمېږي۔ نو په~$\delta_5$
کښې، او هم په $\delta_2$، $\delta_3$ يا~$\delta_4$ کښې،
د نورو د تر ټولو لوړې رتبې پرېکونونو امکان له منځه ځي۔

\end{document}
